% Options for packages loaded elsewhere
\PassOptionsToPackage{unicode}{hyperref}
\PassOptionsToPackage{hyphens}{url}
%
\documentclass[
  11pt,
]{article}
\usepackage{lmodern}
\usepackage{amssymb,amsmath}
\usepackage{ifxetex,ifluatex}
\ifnum 0\ifxetex 1\fi\ifluatex 1\fi=0 % if pdftex
  \usepackage[T1]{fontenc}
  \usepackage[utf8]{inputenc}
  \usepackage{textcomp} % provide euro and other symbols
\else % if luatex or xetex
  \usepackage{unicode-math}
  \defaultfontfeatures{Scale=MatchLowercase}
  \defaultfontfeatures[\rmfamily]{Ligatures=TeX,Scale=1}
\fi
% Use upquote if available, for straight quotes in verbatim environments
\IfFileExists{upquote.sty}{\usepackage{upquote}}{}
\IfFileExists{microtype.sty}{% use microtype if available
  \usepackage[]{microtype}
  \UseMicrotypeSet[protrusion]{basicmath} % disable protrusion for tt fonts
}{}
\makeatletter
\@ifundefined{KOMAClassName}{% if non-KOMA class
  \IfFileExists{parskip.sty}{%
    \usepackage{parskip}
  }{% else
    \setlength{\parindent}{0pt}
    \setlength{\parskip}{6pt plus 2pt minus 1pt}}
}{% if KOMA class
  \KOMAoptions{parskip=half}}
\makeatother
\usepackage{xcolor}
\IfFileExists{xurl.sty}{\usepackage{xurl}}{} % add URL line breaks if available
\IfFileExists{bookmark.sty}{\usepackage{bookmark}}{\usepackage{hyperref}}
\hypersetup{
  hidelinks,
  pdfcreator={LaTeX via pandoc}}
\urlstyle{same} % disable monospaced font for URLs
\usepackage[margin=1in]{geometry}
\setlength{\emergencystretch}{3em} % prevent overfull lines
\providecommand{\tightlist}{%
  \setlength{\itemsep}{0pt}\setlength{\parskip}{0pt}}
\setcounter{secnumdepth}{-\maxdimen} % remove section numbering
\usepackage{mathrsfs}
\usepackage{esint}
\setmathfont{Asana Math}
\setlength{\emergencystretch}{3em}

\author{}
\date{23 September 2026}

\begin{document}

\hypertarget{a-uniform-ricci-energy-gap-above-the-classical-einstein-envelope}{%
\section{A uniform Ricci-energy gap above the classical Einstein
envelope}\label{a-uniform-ricci-energy-gap-above-the-classical-einstein-envelope}}

23 September 2026. This is an application of the parameterized
functional theorem, with actual polarized monopole classes. The
elementary energy envelope below follows from classical inequalities.
The added conclusion is a uniform positive separation from that
envelope, persisting under integral Einstein error. No new topological
nonexistence theorem is asserted.

\hypertarget{statement-on-a-realized-family}{%
\subsection{1. Statement on a realized
family}\label{statement-on-a-realized-family}}

Let \[
 M=X\#k\overline{\mathbb {CP}}^{\,2}\#m(S^1\times S^3),\qquad
 A=c_1^2(X)>0,\quad h=\frac{k+4m}{A},\quad 0<h<\frac13,
 \tag{1}
\] where \(X\) is a minimal complex surface of general type. Put \[
 R=\int_M|\operatorname{Ric}|^2,\qquad
 e=\frac{R}{8\pi^2A},\qquad
 F(h)=\frac{(4-\sqrt{1-3h})^2}{9},
\] \[
 \mathcal E(g)=\inf_{\lambda\ne0}
 \frac{\int_M|\operatorname{Ric}-\lambda g|^2}
      {\lambda^2\operatorname{Vol}(g)}.
 \tag{2}
\] Both \(e\) and \(\mathcal E\) are invariant under constant rescaling.

\textbf{Theorem.} Relative to Theorem 1 of
\href{line1_functional_pencil.md}{the functional-pencil proof}, there is
a universal \(\varepsilon>0\), independent of \(X,k,m,h,g\), such that
\[
 \boxed{\mathcal E(g)+(e-F(h))_+\ge\varepsilon h^2.}
 \tag{3}
\] Thus, for example, \(\mathcal E(g)\le\varepsilon h^2/2\) implies
\(e\ge F(h)+\varepsilon h^2/2\). Every Einstein metric on \(M\), if one
exists, necessarily satisfies \[
 \boxed{\frac{R}{8\pi^2A}\ge F(h)+\varepsilon h^2.}
 \tag{4}
\] An actual nonzero monopole projection excludes Einstein constants
zero and positive, so (4) concerns all Einstein candidates, not just a
chosen sign.

The same theorem holds whenever \(T=2\chi+3\tau=A(1-h)\) and each metric
admits an actual conformally solvable Spin\(^c\) structure with
\(p=(a_g^+)^2\ge A\) and \(a^2\ge T\). No regularity of the class
selection in \(g\) is needed. In (1), the polarized sign selection gives
\[
 p\ge A,\qquad a^2=A-k\ge A-k-4m=T.
 \tag{5}
\] This is gauge-theoretic input, not just an intersection-form
calculation; the conformal class keeps its polarization and its selected
structure. For \(b_+=1\), the small-perturbation chamber and, for
handles, the loop insertions are those in the
\href{lebrun_stability_closure.md}{polarized closure}. The source is
\href{https://arxiv.org/pdf/math/0003068}{LeBrun, Proposition 3.2 and
its proof}.

For a concrete family without chamber issues, take
\(X=\Sigma_g\times\Sigma_q\), \(g,q\ge2\). Then \(A=8(g-1)(q-1)\) and
\(b_+>1\). In particular,
\((\Sigma_2\times\Sigma_2)\#2\overline{\mathbb {CP}}^{\,2}\) has
\(A=8,h=1/4\). Formula (3) applies to every smooth metric on this actual
closed manifold. Any Einstein metric on it would satisfy \[
 \frac{R}{64\pi^2}\ge\frac{49}{36}+\frac{\varepsilon}{16}.
\] We do not assume or produce an Einstein metric on it.

\hypertarget{the-classical-envelope-is-not-the-new-result}{%
\subsection{2. The classical envelope is not the new
result}\label{the-classical-envelope-is-not-the-new-result}}

For an Einstein metric, set \[
 x^2=\frac{\int s^2}{8\pi^2A},\qquad
 y^2=\frac{6\int|W^+|^2}{8\pi^2A}.
\] The classical scalar and scalar--Weyl inequalities and Chern--Weil
give \[
 x\ge2,\qquad x+y\ge3,\qquad x^2+8y^2=12(1-h),\qquad e=x^2/4.
 \tag{6}
\] These imply \(x\ge(8-2\sqrt{1-3h})/3\), hence \(e\ge F(h)\). Equality
is feasible in this finite relaxation; it is not claimed to be realized
by a metric. The constraints are the scalar estimate,
\href{https://arxiv.org/pdf/math/0003068}{LeBrun's Theorem 2.3}, and the
exact Chern--Weil identity. This optimization is entirely classical.

The all-metrics infimum in (1) is \(\inf e=1+h\), by
\href{https://arxiv.org/pdf/math/0003068}{LeBrun's Theorem 4.3}. One has
the exact comparison \[
 F(h)-(1+h)=\frac{4h^2}{(1+\sqrt{1-3h})^2}>0.
 \tag{7}
\] The new term in (4) is therefore an additional uniform energy
separation above the already optimized classical restriction, not the
separation in (7). Classical equality rigidity already excludes exact
Einstein saturation away from the Kähler--Einstein endpoint; it does not
alone give a uniform gap for varying manifolds or almost-Einstein
metrics.

\hypertarget{exact-chernweil-reduction-and-the-proof}{%
\subsection{3. Exact Chern--Weil reduction and the
proof}\label{exact-chernweil-reduction-and-the-proof}}

Use tensor Ricci norms and Hilbert--Schmidt Weyl operator norms. Write
\[
 P=\int|\operatorname{Ric}^{\circ}|^2,\quad S_2=\int s^2,\quad
 W_2=\int|W^+|^2.
\] The identities \[
 S_2=4(R-P),\qquad W_2=2\pi^2T-R/12+P/3
\] turn the functional in the companion theorem into \[
 \mathfrak D_\beta
 =K_\beta R+L_\beta P+12\pi^2(\beta+1)T-72\pi^2p,
 \tag{8}
\] where \[
 K_\beta=\frac{(\beta+1)(8-\beta)}{2\beta},\qquad
 L_\beta=\frac{2(\beta+1)(\beta-2)}{\beta}.
\] Let \(z=\sqrt{1-3h}\) and choose \[
 \beta=\beta(h)=\frac{8-2z}{1+2z}\in(2,8).
 \tag{9}
\] Direct substitution gives \[
 K_\beta F(h)+\tfrac32(\beta+1)(1-h)-9=0.
\] Using (5) in (8), with no discarded negative error, yields \[
 \mathfrak D_\beta
 =8\pi^2A K_\beta(e-F(h))+L_\beta P-72\pi^2(p-A)
 \le8\pi^2A K_\beta(e-F(h))+L_\beta P.
 \tag{10}
\]

First fix the Einstein target to \(-g\), and let
\(V=\operatorname{Vol}(g)\), \(r=V^{-1}\int|\operatorname{Ric}+g|^2\).
The functional-pencil theorem, with \(L=8\) and \(a^2\ge T\), supplies a
universal \(c>0\) such that \[
 c(\beta-2)^2\le r+\mathfrak D_\beta/V.
\] Since \(P/V\le r\), equation (10) implies \[
 c(\beta-2)^2
 \le (1+L_\beta)r+
 K_\beta\frac{8\pi^2A}{V}(e-F(h)).
 \tag{11}
\] The scalar monopole bound and the trace inequality give \[
 32\pi^2A\le S_2,\qquad
 \sqrt{S_2/V}\le4+2\sqrt r,
 \qquad \frac{8\pi^2A}{V}\le(2+\sqrt r)^2.
 \tag{12}
\] For \(r\le1\), replace the last factor in (11) by its positive part
before using (12). On \([2,8]\), \[
 0\le K_\beta\le9/2,\qquad 0\le L_\beta\le27/2.
\] Consequently \[
 c(\beta-2)^2\le\frac{81}{2}\bigl[r+(e-F(h))_+\bigr].
 \tag{13}
\] Finally, \[
 \beta(h)-2
 =\frac{18h}{(1+z)(1+2z)}\ge3h.
\] Thus (3), with \(r\) in place of \(\mathcal E\), holds for \(r\le1\)
with coefficient \(2c/9\). For \(r>1\) it is immediate after choosing
\(\varepsilon=\min\{2c/9,1\}\), because \(h^2<1/9\). For any
\(\lambda<0\), apply this result to \((-\lambda)g\): its normalized
error is the corresponding quotient in (2), while \(e\) and \(h\) are
unchanged. Taking the infimum proves (3) first for \(\mathcal E_-\), the
infimum restricted to negative \(\lambda\).

To include positive target constants, put \(C=8\pi^2A\), \(I=\int s\),
and let \(\mathcal E_+\) be the infimum for \(\lambda>0\). The signed
scalar estimate is \[
 \int s_-^2\ge32\pi^2p\ge4C,\qquad s_-:=\max\{-s,0\}.
 \tag{14}
\] For clarity, it follows directly from an ordinary monopole solution:
with \(\rho=|\Phi|^2\), the integrated spinor identity gives
\(\int\rho^2\le\int s_-\rho\), while Hodge projection gives
\(\int\rho^2\ge32\pi^2p\); apply Cauchy--Schwarz. If \(I>0\), then \[
 \frac{I^2}{V}\le\int s_+^2\le S_2-4C\le4R-4C,
 \qquad \mathcal E_+=4-\frac{I^2}{RV}\ge\frac4e.
\] If \(I\le0\), direct minimization gives \(\mathcal E_+=4\). When
\(e\le F(h)+1\le25/9\), therefore, \(\mathcal E_+\ge36/25>1\). When
\(e>F(h)+1\), the energy excess alone exceeds one. Hence
\(\mathcal E_++(e-F(h))_+\ge1\ge\varepsilon h^2\). Together with
\(\mathcal E=\min\{\mathcal E_-,\mathcal E_+\}\), this proves (3) for
all target signs.

Only the functional-pencil theorem is the nonclassical analytic input.
There is no geometric compactness, scalar band, auxiliary-form
assumption, or additional square-cost budget in this transfer.

\hypertarget{why-the-parameter-eight-does-not-give-this-energy-separation}{%
\subsection{4. Why the parameter eight does not give this energy
separation}\label{why-the-parameter-eight-does-not-give-this-energy-separation}}

At \(\beta=8\), \(K_8=0\). Equation (8) loses \(R\) altogether: \[
 \mathfrak D_8=108\pi^2(T-2p/3)+27P/2.
 \tag{15}
\] This is the old topological gap. Varying \(\beta\) retains the
independent energy coordinate and gives the separation in (3).

There is an exact algebraic non-implication check. At \(h=1/4\), set
\(p=A\), \(a^2=T=3A/4\), \(P=0\), \(s=-4\), and prescribe the constant
Weyl norm \(w=\sqrt6|W^+|=8/7\). Choose \[
 V=2\pi^2A(49/36),\qquad R=4V.
\] These moment data satisfy Chern--Weil, the scalar inequality, every
classical bound \(\mathfrak D_\gamma\ge0\), \(\gamma>0\), and
\(e=F(1/4)=49/36\). However, \[
 \mathfrak D_{7/2}=0,\qquad \mathfrak D_8/V=162/49.
 \tag{16}
\] They violate the new parameter-\(7/2\) estimate. They do not violate
any possible universal coefficient in the global parameter-eight member
\(r+\mathfrak D_8/V\ge36c_8\): necessarily \(c_8\le1/11\). Indeed, for
an actual negative Kähler--Einstein metric \(g_0\), the rescalings
\(g=tg_0\) give \[
 r+\mathfrak D_8/V=4(1-1/t)^2+18/t^2
 \ge36/11,
\] with equality at \(t=11/2\). Yet \(162/49>36/11\). This proves
independence from the parameter-eight numerical inequality even at the
largest coefficient it could possibly have.

The data in (16) are a formal curvature-moment model, not a metric
construction. The manifold and its actual monopole-class selection in
Section 1 are realized; realizing the forbidden moments would contradict
the theorem. This argument does not claim independence from every
possible use of the old differential proof machinery.

\hypertarget{the-remaining-topological-input-is-precise}{%
\subsection{5. The remaining topological input is
precise}\label{the-remaining-topological-input-is-precise}}

Equation (3) does not exclude all energies on any new smooth topology.
For normalized Einstein metrics, the missing independent number is \[
 e=\frac{R}{8\pi^2A}=\frac{V}{2\pi^2A}.
 \tag{17}
\] It is not fixed by \(T\), \(A\), or \(p\). Chern--Weil fixes a linear
combination of scalar and Weyl energies, not this scalar/volume
coordinate. An independently proved upper bound
\(e<F(h)+\varepsilon h^2\) for every Einstein candidate would convert
(4) into a new topological obstruction; no such bound is supplied here.
An all-metrics energy \emph{infimum} cannot serve as that upper bound.

For the pointwise negative Ricci input
\(\delta_-=\|\operatorname{Ric}^{\sharp}+I\|_{\infty,\mathrm{op}}\), (3)
implies the raw consequence \[
 4\delta_-^2+(e-F(h))_+\ge\varepsilon h^2.
 \tag{18}
\] Thus a construction with controlled energy below the strengthened
envelope cannot be arbitrarily Ricci-pinched. This is a restriction on
such input metrics, not a proof that their smooth manifold has no
Einstein metric at a different energy.

\texttt{check\_line1\_energy\_envelope\_gap.py} checks the exact
identities, envelope optimization, constants, scale transformation, and
formal non-implication witness. These finite checks do not certify
Seiberg--Witten solvability or the inherited differential theorem. No
priority claim over unexamined later literature is made.

\end{document}
